\section{Method}
\label{sec:method}

\subsection{Preliminaries}

We describe ReSync on top of a diffusion-transformer backbone trained with a standard diffusion or flow-matching objective. For concreteness, we use the stochastic-interpolant formulation of SiT~\citep{ma2024sit}. Given data $x_0$, noise $\epsilon$, and timestep $t$, the model observes
\begin{equation}
x_t=\alpha_t x_0+\sigma_t\epsilon .
\label{eq:stochastic_interpolant}
\end{equation}
Let $u_t=\dot{\alpha}_t x_0+\dot{\sigma}_t\epsilon$ be the corresponding velocity target. The backbone $f_\zeta$ and output head $v_\zeta$ are trained with
\begin{equation}
\mathcal{L}_{\mathrm{gen}}
=
\mathbb{E}_{x_0,\epsilon,t}
\left[
\left\|v_\zeta(x_t,t,c)-u_t\right\|_2^2
\right].
\label{eq:flow_matching_loss}
\end{equation}
During the same forward pass, the transformer exposes hidden states $h_0(t),\ldots,h_L(t)$, where $h_\ell(t)$ denotes the representation after block $\ell$. ReSync adds training-time losses on these states. The sampling network is unchanged unless the optional layer-skipping branch is enabled.

\subsection{ReSync: Learning Internal Transitions}

Representation alignment methods usually choose a source representation and a target representation, optionally insert a projection head, and then minimize a similarity loss. REPA uses external targets, while SRA and LayerSync show that useful targets can also come from the diffusion model itself~\citep{yu2024repa,jiang2025sra,haghighi2026layersync}. Projection heads can be viewed as practical adapters for matching different representation spaces, but when an adapter is tied to a fixed pair it remains a local solution.

ReSync turns this local adapter into a shared transition model. For two layers $a<b$, we predict the later hidden state from the earlier one as
\begin{equation}
\hat h_b(t)=T_\psi(h_a(t);a,b,t),
\label{eq:resync_predictor}
\end{equation}
where $T_\psi$ is a lightweight predictor conditioned on the source layer, target layer, and timestep. The same predictor is reused across sampled layer pairs, so the auxiliary task is not simply to fit one projection head for one target. Instead, it asks the model to learn reusable changes between internal representation spaces. A cross-timestep variant can replace the target $h_b(t)$ with a cleaner or noisier state from the same sample, following the spirit of SRA; we keep the same-timestep form as the main objective for clarity and lower training cost.

\subsection{Grounded and Compositional Objectives}

The grounded transition loss uses the real target hidden state as a detached supervision signal:
\begin{equation}
\mathcal{L}_{\mathrm{align}}
=
\mathbb{E}_{a<b,t}
\left[
d\left(T_\psi(h_a(t);a,b,t),\mathrm{sg}(h_b(t))\right)
\right],
\label{eq:align_loss}
\end{equation}
where $\mathrm{sg}(\cdot)$ denotes stop-gradient and $d$ is a tokenwise distance, cosine distance by default. This term provides the same kind of internal guidance as self-alignment, but the source layer is not forced to match the target layer in raw feature coordinates.

Grounded prediction alone could still be solved by pair-specific shortcuts. ReSync therefore requires transitions to compose across depth. For layers $a<b<c$, the transition obtained by going through $b$ should agree with the direct transition from $a$ to $c$. We implement this with an EMA copy $\overline{T}_\psi$ as the bootstrap target:
\begin{equation}
\mathcal{L}_{\mathrm{cons}}
=
\mathbb{E}_{a<b<c,t}
\left[
d\left(
T_\psi(T_\psi(h_a(t);a,b,t);b,c,t),
\mathrm{sg}\left(\overline{T}_\psi(h_a(t);a,c,t)\right)
\right)
\right].
\label{eq:consistency_loss}
\end{equation}
This consistency term is the key distinction from a collection of projection heads. It couples short and long transitions, encouraging a shared internal dynamics rather than independent alignment rules for isolated layer pairs.

\subsection{Optional Output Distillation and Training Objective}

ReSync can also be used to test layer-skipping behavior. Given a predicted state $\hat h_b(t)$, we resume the remaining blocks and output head $F_{b+1:L}$:
\begin{equation}
\hat y_{\mathrm{skip}}
=
F_{b+1:L}(\hat h_b(t),t,c).
\label{eq:skip_output}
\end{equation}
The skipped output is matched to the full model output,
\begin{equation}
\mathcal{L}_{\mathrm{out}}
=
\left\|
\hat y_{\mathrm{skip}}
-
\mathrm{sg}(\hat y_{\mathrm{full}})
\right\|_2^2 .
\label{eq:output_distillation_loss}
\end{equation}
We treat this term as optional; the main claim of ReSync is representation regularization through transition prediction and compositional consistency.

The final objective is
\begin{equation}
\mathcal{L}
=
\mathcal{L}_{\mathrm{gen}}
+
\lambda_{\mathrm{align}}\mathcal{L}_{\mathrm{align}}
+
\lambda_{\mathrm{cons}}\mathcal{L}_{\mathrm{cons}}
+
\lambda_{\mathrm{out}}\mathcal{L}_{\mathrm{out}} .
\label{eq:total_loss}
\end{equation}
When $\lambda_{\mathrm{out}}=0$, ReSync is purely a training-time self-distillation objective and adds no inference-time computation.
